% Chapter 2: Literature Review
\chapter{Literature Review}\label{ch:literature}

\guide{Cover the concepts, theories and techniques the reader needs to
understand the rest of the thesis. Organise the chapter by topic, cite the
original sources, and keep it focused on what your work builds on.}

This chapter introduces the concepts on which the thesis builds. It begins with
the study of visual attention and the notion of a saliency map
(\cref{sec:overview}), follows the development of computational saliency models
from hand-crafted features to deep networks trained on large datasets
(\cref{sec:concept}), and then reviews how video representations are learned,
with emphasis on the self-supervised video encoders that this thesis uses
(\cref{sec:video-repr}). The last two sections describe the eye-tracking
datasets and the evaluation metrics of the field (\cref{sec:datasets,sec:metrics}).
Specific video saliency models are reviewed in \cref{ch:related}.

%----------------------------------------------------------------------------
\section{Visual Attention and Saliency}\label{sec:overview}

\subsection{Bottom-Up and Top-Down Attention}

The human visual system cannot process the whole visual field at full
resolution. High acuity is limited to the fovea, so the eyes move several times
per second to bring regions of interest into central vision, and the sequence of
fixations reveals what the observer attends to. Two broad influences on this
selection are usually distinguished: \emph{bottom-up} factors, which depend on
the properties of the image, such as local contrast in colour, orientation or
motion; and \emph{top-down} factors, which depend on the observer's goals,
knowledge and task~\cite{itti2001computational,borji2013state}.

The feature-integration theory of Treisman and Gelade~\cite{treisman1980feature}
proposed that simple features are registered in parallel across the visual field,
whereas attention is required to bind them into objects. Building on this view,
Koch and Ullman~\cite{koch1985shifts} introduced the \emph{saliency map}: a
topographic map that combines the conspicuity of locations across feature
dimensions, from which a winner-take-all mechanism selects the next location to
attend, while inhibition of return prevents attention from returning
immediately to the same place. Itti, Koch and Niebur~\cite{itti1998model} turned
this idea into an influential computational model, in which centre--surround
differences of intensity, colour and orientation at several scales are combined
into a single saliency map. Their model, and the many that followed it, defined
saliency in terms of local feature contrast~\cite{borji2013state}. Graph-based
visual saliency, for example, computes activation maps from feature maps and
normalises them through Markov chains~\cite{harel2006graph}, while the SUN model
defines bottom-up saliency as the self-information of features relative to the
statistics of natural images~\cite{zhang2008sun}.

Top-down influences are harder to model because they depend on what the
observer is trying to do. Judd et al.~\cite{judd2009learning} showed that a
model learned from eye-tracking data, combining low-level features with
higher-level detectors for faces and people and with a centre prior, predicted
fixations better than purely bottom-up models. This shift, from
designing saliency features to learning them from human data, set the
direction for the deep models that followed.

\subsection{Attention in Dynamic Scenes}

Video adds time, and with it motion, events and narrative. In dynamic scenes,
gaze is strongly attracted by motion: in the DIEM recordings, the clustering of
gaze across viewers was predicted by motion~\cite{mital2011clustering}. Carmi and
Itti~\cite{carmi2006visual} distinguished features that cause attentional
selection from features that merely correlate with it, using video clips with
abrupt scene cuts, after which the influence of memory and expectation on gaze
is briefly disrupted.

Two further biases are prominent in video. The first is the \emph{centre bias}.
Observers tend to fixate the centre of the screen, and Tatler~\cite{tatler2007central}
showed that this bias cannot be explained by motor tendencies or by the
distribution of image features alone. For dynamic scenes, Tseng et
al.~\cite{tseng2009quantifying} found that the centre bias correlates strongly
with \emph{photographer bias}, the tendency to frame interesting content
centrally, and is influenced by viewing strategy at scene onsets. The second is
the strong attraction of \emph{faces}, particularly of people who speak: in
recordings of conversations, observers look mostly at faces, and especially at
talking faces, and the original soundtrack makes them follow the turn-taking of
the speakers more closely~\cite{coutrot2014how}. Sound thus influences gaze even
when it carries no spatial information~\cite{coutrot2012influence}.

\subsection{Surprise and Prediction}

A distinct account holds that attention is drawn to events that violate
expectations. Itti and Baldi~\cite{itti2009bayesian} defined \emph{Bayesian
surprise} as the difference between an observer's posterior and prior beliefs
after an observation. In their experiments, 72\,\% of the gaze shifts of
observers watching television were directed towards locations more surprising
than average, rising to 84\,\% for regions selected by all observers. Their
model computed surprise from simple simulated early visual neurons, and earlier
work had used it to detect surprising events in video~\cite{itti2005principled}.
The account connects to predictive-coding theories of the visual cortex, in
which neural responses signal the error between the input and a prediction
formed from context~\cite{rao1999predictive}. Whether the prediction errors of
modern learned video models play a similar role is one of the questions
examined in this thesis.

%----------------------------------------------------------------------------
\section{Deep Saliency Prediction for Images}\label{sec:concept}

The arrival of large eye-tracking and attention datasets enabled saliency models
to be learned end to end. SALICON~\cite{jiang2015salicon} collected
mouse-contingent attention data, a scalable proxy for eye tracking, on images
from MS COCO, and thereby provided a dataset large enough to train deep networks
for saliency. Pan et al.~\cite{pan2016shallow} trained convolutional
networks directly on such data, while Huang et al.~\cite{huang2015salicon}
fine-tuned networks pretrained for object recognition and argued that their
semantic features close much of the gap between low-level saliency and human
fixations. Cornia et al.~\cite{cornia2018predicting} refined the prediction
iteratively with an attentive convolutional LSTM and learned a set of Gaussian
prior maps to model the centre bias, a design that recurs in later models,
including the one proposed in this thesis.

The DeepGaze line of work clarified \emph{why} deep models succeed. K\"ummerer
et al.~\cite{kummerer2017understanding} separated the contributions of low- and
high-level features to fixation prediction, and later showed that transfer
learning from pretrained backbones had been the key driver of progress, with a
twist: replacing the VGG19 backbone of DeepGaze~II by ResNet50 features raised
performance considerably, but better ImageNet classifiers, such as EfficientNet,
brought no further gain~\cite{linardos2021deepgaze}. The quality of the
pretrained representation for saliency is therefore not simply its quality for
classification. Borji~\cite{borji2021saliency} reviewed the deep-learning era of
saliency prediction for images and videos, and identified the remaining gap to
human performance, the failure cases of current models and the difficulty of
fair model comparison as open issues.

%----------------------------------------------------------------------------
\section{Learning Representations of Video}\label{sec:video-repr}

\subsection{Spatio-Temporal Convolutional Networks}

Early deep video models extended image convolutions into time. C3D learned
spatio-temporal features with three-dimensional convolutions on video
clips~\cite{tran2015learning}. Carreira and Zisserman inflated two-dimensional
networks pretrained on ImageNet into three dimensions (I3D) and introduced
Kinetics, a large dataset of human actions that became the standard source of
video pre-training~\cite{carreira2017quo}. Xie et al.~\cite{xie2018rethinking}
showed that many of the costly three-dimensional convolutions can be replaced
by separable spatial and temporal ones, producing the S3D network that many
video saliency models of the following years use as their encoder. Tran et
al.~\cite{tran2018closer} analysed such factorisations systematically; their
R(2+1)D block, a spatial convolution followed by a temporal one, is the basis of
the decoder blocks in this thesis.

\subsection{Video Transformers and Dense Prediction}

The Transformer~\cite{vaswani2017attention} replaced recurrence and convolution
by self-attention, which relates every element of a sequence to every other. The
Vision Transformer (ViT)~\cite{dosovitskiy2021image} applied it to images by
splitting them into patches of $16\times16$ pixels that are embedded as tokens.
For video, ViViT~\cite{arnab2021vivit} and TimeSformer~\cite{bertasius2021space}
extract spatio-temporal tokens and factorise attention over space and time to
cope with the long token sequences, while Video Swin~\cite{liu2022video}
restricts attention to local windows, reintroducing an inductive bias of
locality.

A plain ViT produces features at a single resolution, whereas dense prediction
tasks need several. Feature pyramid networks~\cite{lin2017feature} build
multi-scale representations from convolutional backbones with a top-down
pathway and lateral connections. For transformers, DPT~\cite{ranftl2021vision}
assembles tokens from several stages into image-like maps at different
resolutions and fuses them with a convolutional decoder, and Li et
al.~\cite{li2022exploring} showed that a simple pyramid built by up- and
down-sampling the features of a plain ViT suffices for object detection. The
decoder of this thesis follows these designs (\cref{sec:architecture}).

\subsection{Self-Supervised Video Encoders}\label{sec:ssl}

Supervised pre-training ties a representation to a fixed set of labels.
Self-supervised learning instead derives the training signal from the data
itself, which allows much larger and more varied training sets. Masked
autoencoders (MAE)~\cite{he2022masked} mask a large fraction of image patches
and train the network to reconstruct their pixels; VideoMAE~\cite{tong2022videomae}
applied the same principle to video with very high masking ratios.
UMT~\cite{li2023unmasked} instead aligns the unmasked tokens of a video student
with the features of an image foundation model, which speeds up training and
gives semantically rich features; its ViT-L is the encoder of SalFoM
(\cref{sec:related-vfm}).

Joint-embedding predictive architectures (JEPA) take a third route. Rather than
reconstructing pixels, they predict the \emph{representations} of masked parts
of the input. I-JEPA~\cite{assran2023self} predicts the representations of
several target blocks of an image from a single context block, and
V-JEPA~\cite{bardes2024revisiting} extends this to video. With $x$ the visible
(context) part of a clip, $y$ the full clip, $E_\theta$ the encoder,
$\bar E$ an exponential-moving-average copy of it, $P_\phi$ the predictor and
$\mathcal{M}$ the set of masked tokens, the objective can be written
as~\cite{murlabadia2026vjepa21}
\begin{equation}
  \mathcal{L}_{\mathrm{JEPA}} = \frac{1}{|\mathcal{M}|}\sum_{i\in\mathcal{M}}
  \big\| P_\phi\big(E_\theta(x),\Delta_y\big)_i - \mathrm{sg}\big(\bar E(y)_i\big)\big\|_1,
  \label{eq:jepa}
\end{equation}
where $\Delta_y$ are learnable tokens that mark the positions to be predicted and
$\mathrm{sg}$ stops the gradient. Because the target is a representation rather
than pixels, the encoder can ignore unpredictable detail. Bardes et
al.~\cite{bardes2024revisiting} report that the resulting features perform well
on both appearance- and motion-oriented tasks without adapting the encoder,
that is, with a frozen backbone. A related observation for images is that
self-supervised features can serve many tasks without fine-tuning~\cite{oquab2024dinov2}.

Since the loss in \cref{eq:jepa} is applied only to the masked tokens, the
visible tokens receive no direct training signal. Mur-Labadia et
al.~\cite{murlabadia2026vjepa21} found that the features of the V-JEPA~2 models
consequently carry little local structure (22.2 mIoU in linear-probe semantic
segmentation on ADE20K), and proposed V-JEPA~2.1, which adds a loss on the
visible tokens and self-supervision at intermediate layers to obtain denser
features. Both model families are used in this thesis.

A JEPA model also contains its predictor $P_\phi$. After training, the
predictor's error on a new video, the difference between what it predicted for
a region and what the encoder actually computed there, is a learned measure of
how unexpected that region is. This is the quantity examined as a possible
attention cue in \cref{sec:g2}.

\subsection{Adapting Pretrained Models}

A pretrained encoder can be used frozen, with only a task-specific head
trained on top, or adapted to the task. Full fine-tuning updates all parameters
and can overfit small datasets. Parameter-efficient methods adapt a model with
few trainable parameters: adapters insert small bottleneck layers into a frozen
network~\cite{houlsby2019parameter}, and LoRA~\cite{hu2022lora} adds trainable
low-rank updates to the weight matrices of a frozen transformer, matching full
fine-tuning on language tasks while training a small fraction of the
parameters. Conditioning information, such as an audio signal, can be injected
into a network through feature-wise linear modulation (FiLM), which scales and
shifts intermediate features as a function of the conditioning
input~\cite{perez2018film}. Finally, averaging the predictions of several
independently trained models, a deep ensemble, is a simple way of improving
accuracy and calibration~\cite{lakshminarayanan2017simple}. Each of these
techniques is evaluated in this thesis.

%----------------------------------------------------------------------------
\section{Eye-Tracking Datasets for Video}\label{sec:datasets}

\Cref{tab:lit-datasets} lists the video eye-tracking datasets relevant to this
thesis. They differ in size, content and, importantly, viewing task.

\begin{table}[htbp]
  \centering
  \caption{Video eye-tracking datasets relevant to this thesis}\label{tab:lit-datasets}
  \footnotesize
  \begin{tabularx}{\textwidth}{@{}>{\raggedright}p{2.5cm}>{\raggedright}p{1.6cm}>{\raggedright}p{2.4cm}>{\raggedright\arraybackslash}X@{}}
    \toprule
    \textbf{Dataset} & \textbf{Videos} & \textbf{Viewers} & \textbf{Content and task} \\
    \midrule
    DHF1K~\cite{wang2018revisiting,wang2021revisiting} & 1,000 & 17 & varied YouTube videos, free viewing; test labels withheld \\
    LEDOV~\cite{jiang2018deepvs} & 538 & 32 & varied videos, free viewing \\
    Hollywood-2~\cite{marszalek2009actions,mathe2015actions} & 1,707 & 19 & movie clips; most viewers performed action recognition \\
    UCF Sports~\cite{rodriguez2008action,mathe2015actions} & 150 & 19 & sports videos; same protocol as Hollywood-2 \\
    Hollywood-2, Vig et al.~\cite{vig2012space} & 1,707 & 4 + 1 pooled & independent recording during action recognition \\
    DIEM~\cite{mital2011clustering} & 84 & 42 & advertisements, trailers, documentaries; audio on \\
    Coutrot1~\cite{coutrot2012influence,coutrot2014how} & 60 & 72 & four visual categories; with and without sound \\
    Coutrot2~\cite{coutrot2014how} & 15 & 40 & four people in meetings \\
    ETMD~\cite{koutras2015perceptually,tsiami2019behaviorally} & 12 & 10 & clips from Hollywood movies \\
    AVAD~\cite{min2016fixation} & 45 & 16 & short audio-visual scenes \\
    SumMe~\cite{gygli2014creating,tsiami2019behaviorally} & 25 & 10 & user-made videos \\
    \bottomrule
  \end{tabularx}
\end{table}

DHF1K~\cite{wang2018revisiting,wang2021revisiting} is the largest free-viewing
benchmark: 1,000 videos annotated by 17 observers, split into 600 training, 100
validation and 300 test videos whose annotations are withheld and evaluated by
the dataset authors on request. Hollywood-2 and UCF Sports are action
recognition datasets~\cite{marszalek2009actions,rodriguez2008action} for which
Mathe and Sminchisescu~\cite{mathe2015actions} recorded the gaze of 19 observers,
most of whom were asked to recognise the action; their ground truth is therefore
task-driven and concentrated on the actors. Vig et al.~\cite{vig2012space}
independently recorded gaze on Hollywood-2 during action recognition and
released the raw recordings openly. The six audio-visual datasets of the
lower part of \cref{tab:lit-datasets} were assembled into a common benchmark by
Tsiami et al.~\cite{tsiami2020stavis}; the viewer counts in the table follow
their description and that of later work~\cite{xiong2024diffsal}. Fixation maps
in all these datasets are converted into continuous density maps by Gaussian
smoothing, and the binary fixation maps are kept for location-based metrics.

The differences in task are central to interpreting results across datasets.
Free viewing (DHF1K, DIEM) captures where people look spontaneously; the action
recognition task (Hollywood-2, UCF Sports) concentrates gaze on the actors and
their movements; and in conversations (Coutrot2) gaze follows the speaker.
A model trained on one condition need not transfer to another, which is why
results on these datasets are usually obtained after fine-tuning on each.

%----------------------------------------------------------------------------
\section{Evaluation Metrics}\label{sec:metrics}

No single metric captures all aspects of a saliency map. Bylinskii et
al.~\cite{bylinskii2019metrics} analysed the common metrics and grouped them into
\emph{location-based} metrics, which evaluate the predicted map at the fixated
locations, and \emph{distribution-based} metrics, which compare the predicted map
with the continuous ground-truth density. The location-based family includes the
normalised scanpath saliency (NSS)~\cite{peters2005components}, the mean of the
standardised prediction at fixated pixels, and the area under the ROC curve in
the variant of Judd et al.\ (AUC-J)~\cite{judd2009learning}, which separates
fixated from non-fixated pixels. The shuffled AUC (sAUC) draws its negatives
from the fixations made on other images or videos, so that a model is not
rewarded merely for predicting the centre
bias~\cite{zhang2008sun,borji2013quantitative}. The distribution-based family
includes the linear correlation coefficient (CC), the histogram intersection
(SIM) and the Kullback--Leibler divergence (KLD). Le Meur and
Baccino~\cite{lemeur2013methods} discuss the strengths and weaknesses of these
comparisons, and \cref{sec:evaluation-metrics} gives the exact definitions used
in this thesis.

The metrics respond differently to the same map~\cite{bylinskii2019metrics}.
KLD heavily penalises predictions that assign almost no probability to regions
that observers fixate; SIM measures the overlap of two distributions and is
sensitive to ground-truth mass that the prediction misses; the AUC metrics
depend only on the ranking of saliency values; and sAUC penalises centre bias.
K\"ummerer et al.~\cite{kummerer2015information} proposed to compare models in
terms of information gain over a baseline, which unifies several metrics in a
probabilistic framework, and later showed that the ranking of models depends on
how a probabilistic prediction is turned into a map for each metric. They argued
that each metric should be given the map that is optimal for it under the
model's prediction~\cite{kummerer2018saliency}; in practice, video saliency
benchmarks score one map per model with all metrics.
In video saliency practice, the metrics are computed per frame, averaged over
the frames of each video and then over videos, usually with evaluation code
derived from the MIT saliency benchmark~\cite{bylinskii2019metrics}; small
differences between implementations make numbers from different papers only
approximately comparable.

%----------------------------------------------------------------------------
\section{Summary}\label{sec:literature-summary}

Visual attention is shaped by image features, by the observer's task and by
strong regularities such as the centre bias and the attraction of faces, and in
video additionally by motion, events and sound. Computational models moved from
hand-crafted feature contrast to deep networks whose success rests largely on
pretrained representations. Video representation learning followed a parallel
path from three-dimensional convolutions to transformers and, most recently, to
self-supervised encoders trained on large video collections; among these, the
joint-embedding predictive models are distinctive in learning by predicting
representations, and in containing a predictor whose errors could indicate what
is unexpected. The datasets and metrics of the field are diverse in task and in
what they reward, which makes careful and consistent evaluation essential. The
next chapter reviews how these ingredients have been combined into video
saliency models.
